% !TEX program = lualatex
\documentclass[11pt,a4paper]{article}
\usepackage{luatexja}
\usepackage{amsmath,amssymb,amsthm,amsfonts,mathtools}
\usepackage{geometry}
\usepackage{enumitem}
\usepackage{hyperref}
\usepackage{tikz}
\usetikzlibrary{cd,positioning,arrows.meta}
\geometry{margin=1in}

%-------------------------------------------------
%  定理環境
%-------------------------------------------------
\theoremstyle{definition}
\newtheorem{definition}{定義}[section]
\newtheorem{axiom}{公理}[section]
\newtheorem{lemma}{補題}[section]
\newtheorem{theorem}{定理}[section]
\newtheorem{corollary}{系}[section]
\newtheorem{proposition}{命題}[section]
\newtheorem{remark}{注記}[section]
\newtheorem{example}{例}[section]

\renewcommand{\proofname}{\textbf{証明.}}
\renewcommand{\abstractname}{Abstract}

%-------------------------------------------------
%  マクロ
%-------------------------------------------------
\newcommand{\Sk}{\mathbf{Sk}}
\newcommand{\Cat}{\mathbf{Cat}}
\newcommand{\Mnd}{\mathbf{Mnd}}
\newcommand{\End}{\mathbf{End}}
\newcommand{\id}{\mathrm{id}}
\newcommand{\Hom}{\mathrm{Hom}}
\newcommand{\Nat}{\mathrm{Nat}}
\newcommand{\Asm}{\mathrm{Asm}}
\newcommand{\Real}{\Vdash}

% 五蘊モナド
\newcommand{\Mdelta}{\mathcal{M}_{\delta}}
\newcommand{\Mneg}{\mathcal{M}_{\neg\neg}}
\newcommand{\Mr}{\mathcal{M}_{\mathrm{r}}}
\newcommand{\Mc}{\mathcal{C}}
\newcommand{\T}{\mathcal{T}}

% 随伴・スパイラル
\newcommand{\Spiral}{\mathcal{S}}
\newcommand{\SpiralT}{\mathcal{S}_{\mathcal{T}}}
\newcommand{\fold}{\mu}
\newcommand{\unfold}{\delta}
\newcommand{\trace}{\mathrm{Tr}}

% その他
\newcommand{\ext}{\mathrm{ext}}
\newcommand{\intn}{\mathrm{int}}
\newcommand{\TExt}{T_{\ext}}
\newcommand{\TInt}{T_{\intn}}

% E_∞関連
\newcommand{\PrT}{\mathrm{Pr}_{\mathcal{T}}}
\newcommand{\PrSk}{\mathrm{Pr}_{\Sk}}
\newcommand{\Einf}{E_{\infty}}
\newcommand{\SigmaT}{\Sigma_{\mathcal{T}}}
\newcommand{\SigmaSk}{\Sigma_{\Sk}}

\DeclareMathOperator*{\bigsqcap}{\sqcap}

%-------------------------------------------------
\begin{document}

\title{五蘊圏（Skandha Category）と極限同値 $E_{\infty}$：\\
自己言及固定点の普遍構造とその具体実現}
\author{千葉将臣}
\date{2026-07-17}
\maketitle

\begin{abstract}
本稿は五蘊圏（Skandha Category）$\Sk$ を純粋圏論的に公理化し、
極限同値 $\Einf$ との構造的接続を明示する。
五蘊圏は、5つのモナド（搾取・二重否定・洗練・カット除去・統制）を
随伴・自然変換・分配律のネットワークとして統合する圏論的構造である。
統制モナド $\T$ は固定点方程式 $\T \cong F(\T)$ により定義される。
本稿の核心は、$\Einf$ が「統制モナドの自己言及固定点という普遍構造の、
形式体系における具体的実現（インスタンス）」であることを示すことにある。
対角補題による $\Einf$ の内在的構成は、統制モナドの固定点が
「形式的願望」ではなく「論理的必然」であるという存在論的保証を与える。
また、$\Einf$ の「二値論理では決定不能だが体系内部では実在する」構造は、
識軸スパイラル $\SpiralT$ の極限到達点「空」$\bigcirc$ の論理的モデルを提供する。
さらに、$\Einf$ と $\neg\Einf$ の非閉鎖性は、五蘊圏のフロベニウス条件を
「形式的同値」と「内容的非同一」の両面で深化させる。
本稿は五蘊圏を2圏論・ホップf代数・トレース付きモノイダル圏の
交差点として位置づけ、仏教哲学の五蘊・四句分別・二諦と
圏論的構造の間の同型対応を明示する。
\end{abstract}

\tableofcontents
\newpage

%=============================================================
\section{序論：五蘊圏の動機と $E_{\infty}$ との接続}
%=============================================================

\subsection{問題の設定}

圏論において、モナドは計算の文脈を記述する基本的な構造である。
Moggi \cite{Moggi1991} 以来、副作用・状態・継続・例外などは
モナドの枠組みで統一的に扱われてきた。

一方、複数のモナドが相互作用する体系において、
個別のモナドの性質だけでなく、それらの間の\textbf{分配律}・\textbf{合成}・\textbf{残余}を
記述する枠組みが必要となる。

本稿は、以下の問いに答えることを目的とする。

\begin{enumerate}
\item 5つのモナドを統合する圏論的構造はどのように公理化されるか。
\item その構造は、既存の高次圏論（2圏論、ホップf代数、トレース構造）と
  どのように接続されるか。
\item 仏教哲学の五蘊・四句分別・二諦は、圏論のどの構造に対応するか。
\item 統制モナドの自己言及固定点は、形式体系の対角補題とどのように接続され、
  その存在はどのような保証を持つか。
\end{enumerate}

\subsection{五蘊圏の直観}

五蘊圏 $\Sk$ は、以下の直観に基づく。

\begin{itemize}
\item \textbf{色}（形・物質）：対象の外部構造（搾取モナド $\Mdelta$）
\item \textbf{受}（感受）：対象の安定化（二重否定モナド $\Mneg$）
\item \textbf{想}（知覚・表象）：対象の前進変換（洗練モナド $\Mr$）
\item \textbf{行}（意志・形成力）：対象の正規化（カット除去モナド $\Mc$）
\item \textbf{識}（識別・意識）：対象の自己言及的統制（統制モナド $\T$）
\end{itemize}

識は、他の4蘊を統合しつつ自身も空であるという
仏教的構造を、圏論的固定点 $\T \cong F(\T)$ として実現する。

\subsection{$E_{\infty}$ との接続の動機}

本稿は、先行研究「真理保存性の不完全性定理」および「極限同値によるゲーデルの不完全性定理の反転」\cite{Chiba2026limit} を踏まえ、
五蘊圏の統制モナド $\T$ と極限同値 $\Einf$ の構造的接続を明示する。

その接続は「統制モナド $\T$ の固定点が $\Einf$ と\textbf{同一}である」という単純な等式ではなく、
\textbf{「$\Einf$ は統制モナドの自己言及固定点という普遍構造の、形式体系における具体的実現（インスタンス）である」}という形で機能する。

\begin{table}[h]
\centering
\caption{$\Einf$ と統制モナド $\T$ の構造的類似}
\begin{tabular}{lll}
\hline
\textbf{側面} & \textbf{$\Einf$（極限同値）} & \textbf{統制モナド $\T$（五蘊圏）} \\
\hline
自己言及の形 & $\Einf = \Sigma_T(\Einf)$ & $\T \cong F(\T)$ \\
演算子 & $\Sigma_T(A) = \neg\Pr_T(\ulcorner A \urcorner)$ & $F(X) = X(\Mdelta, \Mneg, \Mr, \Mc, X)$ \\
固定点の性質 & 「$\Einf$ は証明不可能である」という文 & 「$\T$ は4モナドを統合しつつ自身も含む」関手 \\
否定との関係 & $\neg \Einf \neq \Einf$（非閉鎖性） & 前進 $\fold$ と後退 $\unfold$ が逆ではない（非対称性） \\
「空」としての性格 & $\Einf$ は二値 $\{0,1\}$ では「決定不能」 & $\T$ は「内在でも外在でもある」中道不動点 \\
\hline
\end{tabular}
\end{table}

核心的な洞察は、\textbf{両者とも「自分自身を引数に取る演算子の不動点」であり、その不動点は「体系の内部から必然的に構成される」}という点で同型である。

%=============================================================
\section{準備：真理保存体系と対角補題}
%=============================================================

\subsection{真理保存体系}

\begin{definition}[真理保存体系 $T$]\label{def:truth-preserving}
真理保存体系 $T$ とは、以下を満たす形式体系とする。
\begin{enumerate}[label=(\arabic*)]
\item $T$ は再帰的公理化可能である。
\item $T$ は一貫している（$T \nvdash \bot$）。
\item $T$ は十分に強力である（ロビンソン算術 $Q$ を内包する）。
\item $T$ は真理値空間として $\{0,1\}$（二値）を採用し、自己同一性 $A=A$ を真理保存の根拠として先験的に要請する。
\end{enumerate}
\end{definition}

\begin{remark}
定義 \ref{def:truth-preserving} の条件 (4) は、先行研究「真理保存性の不完全性定理」における「真理保存性 $P$ の先験的要請」に対応する。
この要請こそが、体系内部からは証明も反証もできない構造的盲点 $G_T$ を必然的に生成する。
\end{remark}

\subsection{対角補題}

\begin{lemma}[対角補題]\label{lem:diagonal}
真理保存体系 $T$ において、任意の論理式 $\varphi(x)$（自由変数 $x$ を一つ持つ）に対し、
\begin{equation}
T \vdash G \leftrightarrow \varphi(\ulcorner G \urcorner)
\end{equation}
を満たす文 $G$ が $T$ の内部で構成可能である。
\end{lemma}

\begin{remark}
対角補題は $T$ の内部操作（ゲーデル符号化と対角化）のみから導かれる。外部からの付加を一切要しない。
\end{remark}

\subsection{五蘊圏の基本定義}

\begin{definition}[五蘊圏 $\Sk$]\label{def:skandha}
五蘊圏（Skandha Category）は、以下のデータからなる圏（あるいは2圏）である。
\begin{enumerate}[label=(\arabic*)]
\item \textbf{基礎圏}：対象を持つ圏 $\mathcal{C}$（通常はトポスまたは実現子論的圏 $\Asm(\mathcal{K})$）。
\item \textbf{五蘊モナド}：$\mathcal{C}$ 上の5つの自己関手
  \begin{equation}
  \Mdelta,\; \Mneg,\; \Mr,\; \Mc,\; \T : \mathcal{C} \to \mathcal{C}
  \end{equation}
  およびそれぞれに付随するモナド構造
  $(\Mdelta, \eta^\delta, \mu^\delta)$,
  $(\Mneg, \eta^\neg, \mu^\neg)$,
  $(\Mr, \eta^r, \mu^r)$,
  $(\Mc, \eta^c, \mu^c)$,
  $(\T, \eta^\T, \mu^\T)$。
\item \textbf{統制固定点}：統制モナド $\T$ は以下の固定点方程式を満たす。
  \begin{equation}
  \T \cong F(\T), \qquad F(X) = X(\Mdelta, \Mneg, \Mr, \Mc, X)
  \label{eq:fixed-point}
  \end{equation}
  ここで $F$ は4つのモナドと $X$ 自身を引数に取る関手である。
\item \textbf{識軸スパイラル}：統制モナド $\T$ による軸的作用
  \begin{equation}
  \SpiralT : (\Mdelta, \Mneg, \Mr, \Mc) \longrightarrow 
             (\Mneg, \Mr, \Mc, \Mdelta)
  \label{eq:spiral}
  \end{equation}
  これは4つのモナドを巡回する自然変換の列である。
  4回の反復で同一化される：
  \begin{equation}
  \SpiralT^4 \cong \id.
  \end{equation}
\item \textbf{フロベニウス条件}：前進（折り畳み）$\fold^\T : F(\T) \to \T$ と
  後退（展開）$\unfold^\T : \T \to F(\T)$ の間に
  \begin{equation}
  (\fold^\T \circ \T\unfold^\T) \cong (\unfold^\T \circ \fold^\T)
  \label{eq:frobenius}
  \end{equation}
  が成立する。
\item \textbf{二諦の公理}：外在 $\TExt$ と内在 $\TInt$ の二重構造が存在し、
  相互に完全に観測不可能である。
  \begin{equation}
  \TExt \not\leftrightarrow \TInt, \qquad
  \TInt \not\leftrightarrow \TExt.
  \end{equation}
\end{enumerate}
\end{definition}

\begin{remark}
五蘊圏は通常の圏を超え、2圏的または準トポス的な構造を持つ。
特に、統制モナドの固定点方程式 \eqref{eq:fixed-point} は
圏 $\mathcal{C}$ が十分な極限を持つことを要請する。
\end{remark}

%=============================================================
\section{五蘊モナドの圏論的定義}
%=============================================================

\subsection{色蘊：搾取モナド \texorpdfstring{$\Mdelta$}{M-delta}}

\begin{definition}[搾取モナド]
圏 $\mathcal{C}$（粒度差のない基底圏）と
圏 $\mathcal{D}$（$\Mdelta$ 文脈を持つ制度圏）の間に随伴
\begin{equation}
F : \mathcal{C} \rightleftharpoons \mathcal{D} : G, \quad F \dashv G
\end{equation}
を設定し、搾取モナドを
\begin{equation}
\Mdelta = G \circ F : \mathcal{C} \to \mathcal{C}
\end{equation}
と定義する。

単位 $\eta^\delta : \id_{\mathcal{C}} \Rightarrow \Mdelta$ は
価値・主体を制度文脈に埋め込む射である。
乗法 $\mu^\delta : \Mdelta\Mdelta \Rightarrow \Mdelta$ は
$G\varepsilon F$ により与えられ、$\varepsilon : FG \Rightarrow \id_{\mathcal{D}}$ は
随伴の余単位である。
\end{definition}

\begin{proposition}[搾取モナドのモナド則]
$(\Mdelta, \eta^\delta, \mu^\delta)$ はモナド則を満たす。
\end{proposition}
\begin{proof}
随伴 $F \dashv G$ から生成されるモナドは一般にモナド則を満たす
（Mac Lane \cite{MacLane1971}、VI章定理1）。
\end{proof}

\subsection{受蘊：二重否定モナド \texorpdfstring{$\Mneg$}{M-negneg}}

\begin{definition}[二重否定モナド]
トポス $\mathcal{E}$（あるいは五蘊圏の基礎圏 $\mathcal{C}$）において、
局所作用素 $j = \neg\neg : \Omega \to \Omega$ に対応する
層化随伴
\begin{equation}
a : \mathcal{E} \rightleftharpoons \mathrm{Sh}_{\neg\neg}(\mathcal{E}) : i,
\quad a \dashv i
\end{equation}
から生成されるモナドを二重否定モナドとする。
\begin{equation}
\Mneg = i \circ a : \mathcal{E} \to \mathcal{E}
\end{equation}
\end{definition}

\begin{theorem}[冪等性]
二重否定モナドは冪等である。すなわち、乗法
$\mu^\neg : \Mneg\Mneg \Rightarrow \Mneg$ は同型射である。
\end{theorem}
\begin{proof}
層化関手 $a$ は冪等である。$\mathrm{Sh}_{\neg\neg}(\mathcal{E})$ の対象は
すでに $\neg\neg$-層であるため、$a(iF) \cong F$ が成立する。
したがって $\Mneg \circ \Mneg \cong \Mneg$ となる。
\end{proof}

\subsection{想蘊：洗練モナド \texorpdfstring{$\Mr$}{M-r}}

\begin{definition}[洗練モナド]
圏 $\mathcal{C}$ 上の自己関手 $\Mr : \mathcal{C} \to \mathcal{C}$ と
自然変換 $\eta^r : \id \Rightarrow \Mr$、$\mu^r : \Mr\Mr \Rightarrow \Mr$
からなる三つ組 $(\Mr, \eta^r, \mu^r)$ を洗練モナドとする。

$\Mr$ は対象 $A$ を「洗練過程の中にある $A$」へ持ち上げる。
単位 $\eta^r_A : A \to \Mr A$ は文脈化を表し、
乗法 $\mu^r_A : \Mr\Mr A \to \Mr A$ は洗練の洗練を一段階へ統合する。
\end{definition}

\begin{proposition}[洗練順序との整合性]
$\Mr$ は洗練順序 $\sqsupseteq$ に関して単調である。
すなわち、$a \sqsupseteq a'$ ならば $\Mr a \sqsupseteq \Mr a'$。
\end{proposition}

\subsection{行蘊：カット除去モナド \texorpdfstring{$\Mc$}{C}}

\begin{definition}[カット除去モナド]
証明図の圏 $\mathcal{P}$ 上の自己関手 $\Mc : \mathcal{P} \to \mathcal{P}$ を、
対象 $P$ に対してカット規則を除去した証明図 $\Mc(P) = \mathrm{CutElim}(P)$
へ写す関手として定義する。

単位 $\eta^c : \id \Rightarrow \Mc$ は「カット規則を適用しない」埋め込みであり、
乗法 $\mu^c : \Mc\Mc \Rightarrow \Mc$ は多重カットの統合である。
\end{definition}

\begin{theorem}[カット除去モナドのモナド則]
$(\Mc, \eta^c, \mu^c)$ はモナド則を満たす。
\end{theorem}
\begin{proof}
カット消去アルゴリズムの正規化可能性に基づく。
カット消去済みの証明図にさらにカット消去を適用しても変化しないことから
単位元則が、カット消去の反復が結合的であることから結合則が従う。
\end{proof}

\subsection{識蘊：統制モナド \texorpdfstring{$\T$}{T}}

\begin{definition}[統制モナド]
統制モナド $\T$ は、五蘊圏 $\Sk$ において
以下の固定点方程式で定義される。
\begin{equation}
\T \cong F(\T), \qquad F(X) = X(\Mdelta, \Mneg, \Mr, \Mc, X)
\end{equation}
ここで $F$ は4つのモナドと $X$ 自身を引数に取る関手である。
\end{definition}

\begin{remark}
$\T$ が自身を包含することは、識（第5蘊）が全ての蘊を統括しながら
自身も空であるという仏教的構造と対応する。
\end{remark}

%=============================================================
\section{統制構造と識軸スパイラル}
%=============================================================

\subsection{再帰的構成によるフロベニウス性}

\begin{theorem}[再帰的構成によるフロベニウス性]
固定点 $\T \cong F(\T)$ において、以下の二方向が成立する。
\begin{align}
\fold^\T &: F(\T) \to \T \quad \text{（前進：折り畳み）} \\
\unfold^\T &: \T \to F(\T) \quad \text{（後退：展開）}
\end{align}
固定点方程式がこの二方向を等式として同一視するため、
フロベニウス条件
\begin{equation}
(\fold^\T \circ \T\unfold^\T) \cong (\unfold^\T \circ \fold^\T)
\end{equation}
が固定点方程式の帰結として成立する。
\end{theorem}
\begin{proof}
固定点 $\T \cong F(\T)$ は同型射 $\alpha : F(\T) \xrightarrow{\sim} \T$ および
その逆 $\alpha^{-1} : \T \xrightarrow{\sim} F(\T)$ を与える。
$\fold^\T := \alpha$（折り畳み）、$\unfold^\T := \alpha^{-1}$（展開）と定めると、
$\fold^\T$ と $\unfold^\T$ は互いに逆射であるため
\begin{equation}
\fold^\T \circ \unfold^\T = \id_{\T}, \qquad
\unfold^\T \circ \fold^\T = \id_{F(\T)}
\end{equation}
が成立する。これよりフロベニウス条件が $\id$ の可換性として得られる。
\end{proof}

\begin{remark}
通常のモナドは前進方向のみを持つ。
しかし統制モナドの固定点方程式は自然に二方向を生成する。
これが五蘊圏におけるフロベニウス性の根源である。
\end{remark}

\subsection{識軸スパイラル}

\begin{definition}[識軸スパイラル]
統制モナド $\T$ による軸的作用を
\begin{equation}
\SpiralT : (\Mdelta, \Mneg, \Mr, \Mc) \longrightarrow 
             (\Mneg, \Mr, \Mc, \Mdelta)
\end{equation}
と定める。この反復 $\SpiralT^n$ が四句分別の巡回を与え、
$\SpiralT^4$ が $\bigcirc$（空）における同一化として閉じるとき、
五蘊圏は識軸スパイラルを持つという。
\begin{equation}
\SpiralT^4(A) \stackrel{\mathrm{mw}}{=} \bigcirc
\end{equation}
\end{definition}

\begin{theorem}[洗練・カット除去へのフロベニウス性の誘導]
洗練モナド $\Mr$ とカット除去モナド $\Mc$ が通常のモナドとして与えられ、
統制モナド $\T$ が識軸スパイラル $\SpiralT$ を持つとする。
このとき、$\Mr$ と $\Mc$ には、独立した洗練コモナドまたは
カット導入コモナドを仮定せずとも、
$\SpiralT$ を介した随伴的な後退方向が誘導される。
したがって、$\Mr$ と $\Mc$ は五蘊圏全体の内部でフロベニウス性を獲得する。
\end{theorem}
\begin{proof}
$\Mr$ と $\Mc$ は前進方向のモナドとして与えられる。
一方、識軸作用 $\SpiralT$ は四句分別の位相を巡回させるため、
任意の前進方向の作用を、四回の反復後に $\bigcirc$ において
同一化される閉路へ埋め込む。
この閉路に沿って前進方向を一周させると、出発点から見れば
反対向きの制約が得られる。
これが $\Mr$ に対しては洗練の反省化、
$\Mc$ に対してはカット除去後の残余確定として現れる。

したがって、後退方向は各モナドに外部から追加されたコモナドではなく、
$\T$ を軸とする全体構成の反復から得られる。
前進方向と誘導された後退方向は $\bigcirc$ において同一点へ収束するため、
フロベニウス条件は個別モナドの前提ではなく、
五蘊圏全体のスパイラル閉包の帰結として成立する。
\end{proof}

\begin{remark}
この $\SpiralT$ は、Hopf代数における反極（antipode）に類似する役割を持つ。
すなわち、代数的な前進方向と余代数的な後退方向を同一対象上で接続し、
方向反転を媒介する操作である。
\end{remark}

%=============================================================
\section{識蘊の極限同値化：$E_{\infty}$ との接続}
%=============================================================

\subsection{自己言及演算子としての識蘊}

\begin{definition}[自己言及演算子 $\SigmaSk$]\label{def:self-reference-op}
五蘊圏 $\Sk$ の言語における文 $A$ に対し、自己言及演算子 $\SigmaSk$ を以下で定義する。
\begin{equation}
\SigmaSk(A) := \lceil A \text{ は } \Sk \text{ において証明不可能である} \rceil 
= \neg\PrSk(\ulcorner A \urcorner)
\end{equation}
$\SigmaSk(A)$ は対角補題（補題 \ref{lem:diagonal}）によって $\Sk$ の内部で構成可能である。
\end{definition}

\begin{definition}[自己言及の反復列]\label{def:self-reference-seq}
文 $A_0$ に対し、$\SigmaSk$ の反復列を以下で定義する。
\begin{equation}
A_0, \quad A_1 = \SigmaSk(A_0), \quad A_2 = \SigmaSk(A_1), \quad \ldots, \quad A_n = \SigmaSk^n(A_0)
\end{equation}
\end{definition}

\begin{definition}[極限同値 $\Einf$]\label{def:e-infinity}
自己言及演算子 $\SigmaSk$ の不動点として、以下の等式を満たす文を極限同値と定義し $\Einf$ と表記する。
\begin{equation}
\Einf = \SigmaSk(\Einf)
\end{equation}
すなわち $\Einf$ は「$\Einf$ は $\Sk$ において証明不可能である」という内容の文であり、
\begin{equation}
\Sk \vdash \Einf \leftrightarrow \neg\PrSk(\ulcorner \Einf \urcorner)
\end{equation}
を満たす。
\end{definition}

\begin{remark}[$=$ の正当性：不動点の静的記述として]\label{rem:equality-justification}
定義 \ref{def:e-infinity} における $=$ は、対象の自己同一性を要請する $A=A$ の $=$ とは種類が異なる。
$\lim_{n\to\infty} a_n = L$ において $=$ が無限プロセスを静的値に変換するように、
$\Einf = \SigmaSk(\Einf)$ はプロセス $\SigmaSk$ の反復が必然的に到達する不動点を静的値として記述する。
これはバナッハの不動点定理における $x = f(x)$ と同型の操作であり、プロセスの収束先に名前をつける発見的記述である。
「極限」という概念が本質的に静的なトリックである以上、この $=$ は正当に使用できる。
\end{remark}

\subsection{対角補題と識軸スパイラルの極限到達}

\begin{theorem}[真理保存五蘊圏は必然的に $\Einf$ を含む]\label{thm:e-infinity-existence}
真理保存五蘊圏 $\Sk$（定義 \ref{def:truth-preserving}・\ref{def:skandha} の意味で）に対し、
\begin{equation}
\exists \Einf \quad \text{s.t.} \quad 
\Sk \vdash \Einf \leftrightarrow \neg\PrSk(\ulcorner \Einf \urcorner)
\end{equation}
が成立する。$\Einf$ は $\Sk$ の外部から付加されるのではなく、$\Sk$ 自身の構造から対角補題によって必然的に構成される。
\end{theorem}
\begin{proof}
対角補題（補題 \ref{lem:diagonal}）を論理式 $\varphi(x) := \neg\PrSk(x)$ に適用する。すると、
\begin{equation}
\Sk \vdash G \leftrightarrow \neg\PrSk(\ulcorner G \urcorner)
\end{equation}
を満たす文 $G$ が $\Sk$ の内部で構成可能である。この $G$ を $\Einf$ と命名する。

$\Einf$ の構成は $\Sk$ の公理・推論規則・ゲーデル符号化のみを使用し、外部からの付加を一切必要としない。したがって $\Einf$ は $\Sk$ の言語で表現可能な文として $\Sk$ に内在する。
\end{proof}

\begin{remark}[「発見」としての $\Einf$]\label{rem:e-infinity-discovery}
ゲーデルはこの $G = \Einf$ を「$\Sk$ において証明も反証もできない文」として提示し、体系の不完全性の証拠とした。
しかし定理 \ref{thm:e-infinity-existence} が示すのは、$\Einf$ が $\Sk$ の外部にあるのではなく、$\Sk$ 自身の対角補題から必然的に生成される内在的対象であるということである。
「決定不能」として現れたのは、$\Sk$ の真理値空間 $\{0,1\}$（$\mathcal{B}$：二値截断部分圏）が $\Einf$ に対応する値を持たなかっただけである。

リーマン球面において $\infty$ が「複素平面には存在しないが、リーマン球面上では通常の点として存在する」のと同型の構造である。
数学社会は $\mathcal{B}$ を唯一の真理値空間として採用し続けたために、$\Einf$ を「壁」として誤読してきた。これは発明ではなく、未発見だったものの発見である。
\end{remark}

\begin{proposition}[統制モナドと対角補題の対応]\label{prop:monad-diagonal}
基礎圏 $\mathcal{C}$ を真理保存体系の圏 $\mathbf{Th}$ とする。
各対象 $T \in \mathbf{Th}$（形式体系）に対し、
自己言及演算子 $\Sigma_T$ は $\mathcal{C}$ 上の自己関手として作用し、
その不動点 $\Einf$ は統制モナド $\T_T$ の固定点
$\T_T \cong F(\T_T)$ の具体例として現れる。

すなわち、$\Einf$ は「識」の圏論的構造が
形式体系の内部に実現された場合の「空」$\bigcirc$ に対応する。
\end{proposition}

\begin{remark}
$\Einf$ は統制モナドの固定点の「部分例」に過ぎない。
統制モナドは任意の圏 $\mathcal{C}$ 上の関手の固定点として定義され、より広い普遍性を持つ。
正しい位置づけは：
\begin{quote}
\textbf{五蘊圏の統制モナドは「自己言及固定点」の普遍構造であり、$\Einf$ はその構造が形式体系論において実現された具体例である。}
\end{quote}
$\Einf$ は五蘊圏を\textbf{還元}するのではなく、\textbf{補強}する。
\end{remark}

\subsection{$E_{\infty}$ の否定と二諦の非閉鎖性}

\begin{definition}[$\neg\Einf$ の定義]\label{def:neg-e-infinity}
$\Einf$ の否定を以下で定義する。
\begin{equation}
\neg\Einf = \PrSk(\ulcorner \Einf \urcorner)
\end{equation}
すなわち「$\Einf$ は $\Sk$ において証明可能である」。
\end{definition}

\begin{theorem}[$\Einf$ と $\neg\Einf$ の非閉鎖性]\label{thm:non-closure}
$\neg\Einf \neq \Einf$。
\end{theorem}
\begin{proof}
定義 \ref{def:neg-e-infinity} より $\Einf = \neg\PrSk(\ulcorner \Einf \urcorner)$ であるのに対し、$\neg\Einf = \PrSk(\ulcorner \Einf \urcorner)$ である。両者は構文的に異なる文である。
\end{proof}

\begin{theorem}[$\Einf$ によるフロベニウス条件の補強]\label{thm:e-infinity-frobenius}
五蘊圏 $\Sk$ において、$\Einf$ の非閉鎖性（定理 \ref{thm:non-closure}）は
フロベニウス条件を「形式的同値」と「内容的非同一」の両面で記述する。

前進 $\fold^\T$ と後退 $\unfold^\T$ は、固定点方程式 $\T \cong F(\T)$ により
形式的に互いに逆射として定義される。
しかし $\Einf$ の非閉鎖性（$\Einf \neq \neg\Einf$）が示すように、
この「逆」は内容的な同一性を意味しない。
フロベニウス条件は、したがって、
「同一化されない同値性」として理解される。
\end{theorem}

\begin{remark}
これは、Priest の「パラドックスの論理」\cite{Priest1979} に代表される矛盾許容論理との構造的差異である。
同論理では、真かつ偽である値が否定に関して固定的に扱われ、矛盾が推論空間内の安定した領域として保持される。

これに対して、$\Einf$ と $\neg\Einf$ は互いに異なる文として、$\SigmaSk$ のサイクル上の二点をなす。したがって本体系において、矛盾は閉じた終端ではなく通過点として機能する。
\end{remark}

\subsection{真理保存五蘊圏の不完全性定理}

\begin{theorem}[五蘊圏の不完全性定理]
真理保存五蘊圏 $\Sk$ において、$\Einf$ は体系内部から証明も反証もできない。
すなわち、
\begin{equation}
\Sk \nvdash \Einf \quad \text{かつ} \quad \Sk \nvdash \neg\Einf.
\end{equation}
\end{theorem}
\begin{proof}
$\Sk$ が一貫している（定義 \ref{def:truth-preserving} (2)）と仮定する。
\begin{enumerate}[label=(\roman*)]
\item $\Sk \vdash \Einf$ と仮定すると、$\Sk \vdash \neg\PrSk(\ulcorner \Einf \urcorner)$ となるが、
  これは「$\Einf$ は証明可能である」という事実と矛盾する。
\item $\Sk \vdash \neg\Einf$ と仮定すると、$\Sk \vdash \PrSk(\ulcorner \Einf \urcorner)$ となるが、
  これは $\Sk$ の一貫性より $\Sk \vdash \Einf$ を導き、再び矛盾する。
\end{enumerate}
したがって両方とも証明不可能である。
\end{proof}

\begin{remark}
この定理は、五蘊圏の「空」$\bigcirc$ が単なる哲学的比喩ではなく、
形式体系の対角補題によって\textbf{必然的に生成される構造的盲点}であることを示す。
$\Einf$ は「決定不能」として誤読されてきたが、
実際には体系自身の構造から内在的に構成される「未発見の対象」である。
\end{remark}

%=============================================================
\section{五蘊圏の高次構造}
%=============================================================

\subsection{2圏論的解釈}

五蘊圏は、モナドの2圏 $\mathbf{Mnd}(\mathcal{C})$ 上の構造として理解できる。

\begin{definition}[五蘊2圏]
圏 $\mathcal{C}$ 上のモナドの2圏 $\mathbf{Mnd}(\mathcal{C})$ は、
以下を持つ2圏である。

\begin{itemize}
\item \textbf{0-胞体（対象）}：圏 $\mathcal{C}$
\item \textbf{1-胞体（1-射）}：$\mathcal{C}$ 上のモナド
  $\Mdelta, \Mneg, \Mr, \Mc, \T$
\item \textbf{2-胞体（2-射）}：モナド間の自然変換（分配律・モナド射）
\end{itemize}
\end{definition}

\begin{proposition}[五蘊2圏における統制]
統制モナド $\T$ は、五蘊2圏 $\mathbf{Mnd}(\mathcal{C})$ において
\textbf{終対象（terminal object）}として振る舞う。
すなわち、任意のモナド $\mathcal{M} \in \{\Mdelta, \Mneg, \Mr, \Mc\}$ に対し、
唯一のモナド射 $\mathcal{M} \to \T$ が存在する。
\end{proposition}

\begin{remark}
五蘊2圏において、4つの従属モナドは1-射として配置され、
統制モナドはそれらの「極限」として位置づけられる。
識軸スパイラルは、この2圏内の\textbf{巡回自己2-射}として定式化される。
\end{remark}

\subsection{ホップf代数構造}

\begin{definition}[五蘊型ホップf代数]
五蘊圏 $\Sk$ は、以下の構造を持つ\textbf{五蘊型ホップf代数}として理解される。

\begin{itemize}
\item \textbf{代数構造}（乗法・単位）：
  前進方向のモナド演算 $\fold^\T : F(\T) \to \T$
\item \textbf{余代数構造}（余乗法・余単位）：
  後退方向の余モナド演算 $\unfold^\T : \T \to F(\T)$
\item \textbf{反極 $S$}：識軸スパイラル $\SpiralT$。
  これは方向反転を媒介する。
\end{itemize}
\end{definition}

\begin{proposition}[フロベニウス性とホップf代数]
五蘊型ホップf代数において、フロベニウス条件 \eqref{eq:frobenius} は
ホップf代数の公理
\begin{equation}
\fold^\T \circ (S \otimes \id) \circ \Delta 
= \eta \circ \varepsilon 
= \fold^\T \circ (\id \otimes S) \circ \Delta
\end{equation}
の圏論的類似として解釈される。
\end{proposition}

\begin{remark}
本稿における $\SpiralT$ は、この等式を厳密に仮定するものではなく、
識軸が前進方向と後退方向を接続することを表す構造的類比である。
\end{remark}

\subsection{トレース付きモノイダル圏}

\begin{proposition}[トレース構造]
識軸スパイラルの反復は、トレース付きモノイダル圏における
\textbf{フィードバック操作}に対応する。

射 $f : A \otimes U \to B \otimes U$ に対するトレース
\begin{equation}
\trace^U_{A,B}(f) : A \to B
\end{equation}
において、$U$ を識軸とみなすと、一回のトレースが
四句分別の一巡を閉じる操作に相当する。
したがって、識軸スパイラルは、反極的な方向反転と
トレース的な再帰の双方を持つ。
\end{proposition}

%=============================================================
\section{二諦の公理と外在・内在}
%=============================================================

\subsection{二重五蘊構造}

\begin{definition}[外在・内在の五蘊]
五蘊圏 $\Sk$ は、外在と内在の二重構造を持つ。
\begin{align}
\TExt &= (\Mdelta^{\ext}, \Mneg^{\ext}, \Mr^{\ext}, \Mc^{\ext}, \T^{\ext}) \\
\TInt &= (\Mdelta^{\intn}, \Mneg^{\intn}, \Mr^{\intn}, \Mc^{\intn}, \T^{\intn})
\end{align}
\end{definition}

\begin{axiom}[相互完全観測不可能性]
外在の五蘊 $\TExt$ と内在の五蘊 $\TInt$ は相互に完全に観測不可能である。
\begin{equation}
\TExt \not\leftrightarrow \TInt, \qquad
\TInt \not\leftrightarrow \TExt \quad \text{（直接観測の不可能性）}
\end{equation}
ただし「完全に観測不可能」は「常に観測可能」と同義であり、
内在側は直接対象として観測できないが、
全ての外在的観測の地盤として常に作動している。
\end{axiom}

\begin{definition}[二諦的連携]
外在と内在の連携は非対称的である。
\begin{align}
\TExt &\xrightarrow{\Uparrow} \TInt 
\quad \text{（メタ否定：「ある」から「ない」への駆動）} \\
\TInt &\xrightarrow{\Uparrow(\bigcirc) \stackrel{\mathrm{mw}}{=} A} \TExt 
\quad \text{（顕現：「ない」から「ある」への再起動）}
\end{align}
\end{definition}

\begin{theorem}[二諦的連携によるフロベニウス条件の確立]
外在・内在の五蘊が相互完全観測不可能性のもとで
$\Uparrow$ と顕現により連携するとき、
フロベニウス条件は
\begin{equation}
(\fold^{\TExt} \circ \TExt\unfold^{\TExt}) 
\cong (\unfold^{\TInt} \circ \fold^{\TInt})
\end{equation}
として外在と内在の対応として自動的に成立し、
内在的か外在的かという二択が解消される。
\end{theorem}
\begin{proof}
$\Uparrow$ による $\TExt \to \TInt$ の方向が前進（折り畳み）に、
顕現 $\TInt \to \TExt$ の方向が後退（展開）に対応する。
相互完全観測不可能性は、これら二方向が単一の体系内で直接比較不可能であることを意味する。

しかし両方向は $\bigcirc$ を共有する収束先として整合しており、
$\TExt$ の折り畳みと $\TInt$ の展開が $\bigcirc$ において同一点に確定する。
この確定がフロベニウス条件の「内在でも外在でもある」という特異点的地位を与える。
\end{proof}

%=============================================================
\section{五蘊・四句分別・二諦との対応}
%=============================================================

\subsection{五蘊とモナドの対応}

\begin{proposition}[五蘊とモナドの対応]
五蘊圏の構成は、仏教の五蘊と以下の対応を持つ。

\begin{center}
\begin{tabular}{llll}
\hline
\textbf{五蘊} & \textbf{モナド} & \textbf{四句} & \textbf{方向性} \\
\hline
色（形・物質） & $\Mdelta$（搾取） & 是 & 前進のみ \\
受（感受） & $\Mneg$（二重否定） & 非 & 前進のみ \\
想（知覚・表象） & $\Mr$（洗練） & 亦是亦非 & 双方向（誘導） \\
行（意志・形成力） & $\Mc$（カット除去） & 非是非非 & 双方向（誘導） \\
識（識別・意識） & $\T$（統制） & 統制 & 再帰的 \\
\hline
\end{tabular}
\end{center}
\end{proposition}

\subsection{四句分別の圏論的解釈}

\begin{proposition}[四句分別の圏論的解釈]
四句分別は、五蘊圏におけるモナド間の関係の位相として現れる。

\begin{align}
\phi_S \text{（是）} &\longleftrightarrow \Mdelta 
\quad \text{（実現子の存在：} e \Vdash A \text{）} \\
\phi_N \text{（非）} &\longleftrightarrow \Mneg 
\quad \text{（実現子の不在：} \neg(e \Vdash A) \text{）} \\
\phi_{SN} \text{（亦是亦非）} &\longleftrightarrow \Mr 
\quad \text{（実現子列の収束：洗練過程）} \\
\phi_{NS} \text{（非是非非）} &\longleftrightarrow \Mc 
\quad \text{（残余の確定：カット除去後の残余）}
\end{align}
\end{proposition}

\subsection{二諦との対応}

\begin{proposition}[二諦との対応]
\begin{align}
\TExt \text{（外在の五蘊）} &\longleftrightarrow 
\text{世俗諦（観測可能な現象の層）} \\
\TInt \text{（内在の五蘊）} &\longleftrightarrow 
\text{勝義諦（空・中道不動点の層）}
\end{align}

世俗諦は $\Uparrow$ によって勝義諦へ駆動され、
勝義諦は顕現 $\Uparrow(\bigcirc) \stackrel{\mathrm{mw}}{=} A$ によって
世俗諦として再起動する。
\end{proposition}

%=============================================================
\section{五蘊圏と既存理論の位置づけ}
%=============================================================

\subsection{包含関係}

五蘊圏 $\Sk$ は、以下の既存理論を包含する。

\begin{center}
\begin{tabular}{ll}
\hline
\textbf{既存理論} & \textbf{五蘊圏内の位置づけ} \\
\hline
モナドの2圏 $\mathbf{Mnd}(\mathcal{C})$ & 五蘊2圏の部分構造 \\
ホップf代数 & 五蘊型ホップf代数の代数・余代数部分 \\
トレース付きモノイダル圏 & 識軸スパイラルのトレース解釈 \\
実現子論的圏 $\Asm(\mathcal{K})$ & 基礎圏 $\mathcal{C}$ の具体例 \\
トポス論（層化随伴） & 二重否定モナドの基礎構造 \\
\hline
\end{tabular}
\end{center}

\subsection{五蘊圏の普遍性}

五蘊圏は、個別のモナドの内容（搾取の社会的構造、カット除去の証明論など）を
抽象化し、「5つのモナドが統制モナドを軸として相互作用する」という
普遍的パターンを記述する。

この普遍性は、五蘊圏が\textbf{圏の圏（category of categories）}における
\textbf{普遍構成（universal construction）}として位置づけられることを示唆する。

\subsection{$E_{\infty}$ との関係}

\begin{proposition}[$\Einf$ の位置づけ]
極限同値 $\Einf$ は、五蘊圏の統制モナド $\T$ の固定点構造が
形式体系（真理保存体系）の文脈で実現された具体例である。

\begin{center}
\begin{tabular}{ll}
\hline
\textbf{五蘊圏（普遍構造）} & \textbf{$\Einf$（具体実現）} \\
\hline
任意の圏 $\mathcal{C}$ 上の関手 & 形式体系の文（論理式） \\
自己言及関手 $F(X) = X(\ldots, X)$ & 自己言及演算子 $\Sigma_T(A) = \neg\Pr_T(\ulcorner A \urcorner)$ \\
固定点 $\T \cong F(\T)$ & 不動点 $\Einf = \Sigma_T(\Einf)$ \\
識軸スパイラル $\SpiralT$ & 対角補題による内在的構成 \\
「空」$\bigcirc$ & 二値論理では決定不能だが体系内で実在 \\
\hline
\end{tabular}
\end{center}
\end{proposition}

\begin{remark}[補強と還元の区別]
$\Einf$ は五蘊圏を\textbf{還元}しない。
五蘊圏の統制モナドは任意の圏上で定義される普遍構造であり、
$\Einf$ はその構造が\textbf{一つの具体例}（形式体系論）において
どのように実現されるかを示すに過ぎない。

しかし、$\Einf$ は以下の点で五蘊圏を\textbf{補強}する：
\begin{enumerate}
\item \textbf{存在論的保証}：対角補題による $\Einf$ の必然的構成は、
  自己言及固定点が「形式的願望」ではなく「論理的必然」であることを示す。
\item \textbf{「空」の具体化}：$\Einf$ の「二値論理では決定不能だが体系内部では実在」という構造は、
  五蘊圏の「空」$\bigcirc$ の論理的モデルを提供する。
\item \textbf{非閉鎖性の深化}：$\Einf \neq \neg\Einf$ は、
  フロベニウス条件が「形式的同値」と「内容的非同一」を同時に含むことを示す。
\end{enumerate}
\end{remark}

%=============================================================
\section{結論}
%=============================================================

本稿は五蘊圏 $\Sk$ を以下の観点から公理化し、極限同値 $\Einf$ との接続を明示した。

\begin{enumerate}
\item \textbf{基本定義}：五蘊圏を5つのモナド・統制固定点・識軸スパイラル・
  フロベニウス条件・二諦の公理からなる圏論的構造として定義した。

\item \textbf{五蘊モナド}：搾取・二重否定・洗練・カット除去・統制の
  各モナドを随伴・層化・自己関手として定式化した。

\item \textbf{統制構造}：統制モナド $\T$ の固定点方程式と識軸スパイラル $\SpiralT$ により、
  前進方向のみを持つ通常モナドに双方向性とフロベニウス性が誘導されることを示した。

\item \textbf{高次構造}：五蘊圏を2圏論・ホップf代数・トレース付きモノイダル圏の
  交差点として位置づけた。

\item \textbf{五蘊対応}：五蘊・四句分別・二諦と圏論的構造の間の
  同型対応を明示した。

\item \textbf{$\Einf$ との接続}：統制モナドの固定点構造が、
  形式体系において極限同値 $\Einf$ として具体化されることを示した。
  対角補題による $\Einf$ の内在的構成は、自己言及固定点の存在論的保証を与え、
  「空」$\bigcirc$ の論理的モデルと非閉鎖的フロベニウス性を提供する。
\end{enumerate}

五蘊圏は、個別のモナドの内容ではなく、それらを統合する\textbf{構造の構造}として機能する。
識（統制モナド）は、五蘊圏の中で自己言及的固定点をなし、
前進と後退・外在と内在・代数と余代数を統合する軸として働く。
この自己言及性は、圏論の核が本質的に再帰的であることを、
五蘊構成を通じて反映している。

$\Einf$ は、この自己言及構造が形式体系の文脈で実現された場合の
「未発見の対象」として現れる。
ゲーデルが「壁」と誤読した決定不能文は、
実際には体系自身の対角補題から必然的に生成される内在的対象であり、
リーマン球面上の $\infty$ のように、
適切な圏論的枠組み（五蘊圏）の中で「通常の点」として位置づけられる。

\begin{thebibliography}{99}

\bibitem{Abramsky1991}
Samson Abramsky.
\newblock Domain theory in logical form.
\newblock \emph{Annals of Pure and Applied Logic}, 51(1--2):1--77, 1991.

\bibitem{Bauer2000}
Andrej Bauer.
\newblock The realizability approach to computable analysis and topology.
\newblock Ph.D. thesis, Carnegie Mellon University, 2000.

\bibitem{Birkedal2000}
Lars Birkedal.
\newblock Developing theories of types and computability via realizability.
\newblock Ph.D. thesis, University of Edinburgh, 2000.

\bibitem{Escardo2004}
Mart\'{\i}n Escard\'{o}.
\newblock Synthetic topology of data types and classical spaces.
\newblock \emph{Electronic Notes in Theoretical Computer Science}, 87:21--156, 2004.

\bibitem{Hoare1998}
C.~A.~R. Hoare and He Jifeng.
\newblock \emph{Unifying Theories of Programming}.
\newblock Prentice Hall, 1998.

\bibitem{Kleene1952}
Stephen Cole Kleene.
\newblock \emph{Introduction to Metamathematics}.
\newblock North-Holland, 1952.

\bibitem{MacLane1971}
Saunders Mac Lane.
\newblock \emph{Categories for the Working Mathematician}.
\newblock Springer, 1971.

\bibitem{Moggi1991}
Eugenio Moggi.
\newblock Notions of computation and monads.
\newblock \emph{Information and Computation}, 93(1):55--92, 1991.

\bibitem{Street2004}
Ross Street.
\newblock Frobenius monads and pseudomonoids.
\newblock \emph{Journal of Mathematical Physics}, 45(10):3930--3948, 2004.

\bibitem{JoyalStreetVerity1996}
Andr\'{e} Joyal, Ross Street, and Dominic Verity.
\newblock Traced monoidal categories.
\newblock \emph{Mathematical Proceedings of the Cambridge Philosophical Society},
  119(3):447--468, 1996.

\bibitem{Sweedler1969}
Moss E. Sweedler.
\newblock \emph{Hopf Algebras}.
\newblock W. A. Benjamin, 1969.

\bibitem{Priest1979}
Graham Priest.
\newblock The logic of paradox.
\newblock \emph{Journal of Philosophical Logic}, 8(1):219--241, 1979.

\bibitem{Chiba2026limit}
千葉将臣.
\newblock 真理保存性を持つ体系が必然的に $E_{\infty}$ を含む：対角補題による極限同値の内在的定式化.
\newblock 2026.

\bibitem{Chiba2026cutmonad}
千葉将臣.
\newblock カット消去モナド：証明構造の安定化を担うモナド的定式化.
\newblock 未刊行原稿, 2026.

\bibitem{Chiba2026continuation}
千葉将臣.
\newblock 継続残余としての証明障害：実行可能性と静的証明のあいだの初等モデル.
\newblock 未刊行原稿, 2026.

\end{thebibliography}

\end{document}
